% ---
% id: reductions/rem:branch-blocks
% kind: rem
% title: Conditionally characterized receivers under the branching rule
% summary: >-
%   The branching rule squashes a conditionally characterized receiver as soon
%   as each branch squashes, stacking the branch targets as blocks. A branch
%   that cannot be modelled gets a block of its own, where the adversary fixes
%   the outcome only in proportion to the branch's probability.
% prominence: standard
% status: accepted
% uses: [lem:master]
% assumes: []
% refs: []
% provenance: newly-derived
% lean:
% ---

\begin{remark}\label{rem:branch-blocks}
The branching rule of Lemma~\ref{lem:master} turns any conditional
description of a receiver into a squashing relation. The device is squashable
as soon as \emph{each branch} is, and the squashed target is the branch
targets stacked as diagonal blocks. A branch one cannot model is not fatal: it
gets a block of its own, on which the squashed adversary rather than Bob's
detector fixes the outcome, but only in proportion to how likely that branch
was.
\end{remark}
